\documentclass[11pt]{article}
\usepackage[margin=1in]{geometry}
\usepackage{amsmath,amssymb}
\usepackage{enumitem}
\usepackage{hyperref}
\usepackage{listings}
\usepackage{xcolor}
\usepackage{booktabs}

\lstset{
  basicstyle=\ttfamily\small,
  backgroundcolor=\color{gray!10},
  frame=single,
  framerule=0pt,
  breaklines=true,
  language=R,
  showstringspaces=false
}

\usepackage[many]{tcolorbox}
\definecolor{main}{HTML}{5989cf}
\definecolor{sub}{HTML}{cde4ff}
\definecolor{sub2}{HTML}{fde9ce}

\tcbset{
    sharp corners,
    colback = white,
    before skip = 0.2cm,
    after skip = 0.5cm
}

\newtcolorbox{boxK}{
    sharpish corners,
    boxrule = 0pt,
    toprule = 4.5pt,
    enhanced,
    fuzzy shadow = {0pt}{-2pt}{-0.5pt}{0.5pt}{black!35}
}

\hypersetup{colorlinks=true, linkcolor=blue, urlcolor=blue, citecolor=blue}

\title{Problem Set 6: Event Studies and Synthetic Methods}
\author{MGMT 737 --- Spring 2026}
\date{}

\begin{document}
\maketitle

This problem set covers event study methods (Sun and Abraham, 2021) and synthetic control approaches. You will work with three datasets: the Sun and Abraham (2021) replication data (\texttt{hrs\_eventdata.csv}) from Dobkin et al.\ (2020), the California Proposition 99 cigarette data (built into the \texttt{synthdid} package), and \texttt{wolfers\_aer\_clean.csv} for the staggered synthetic control analysis.

You may use AI coding assistants, but you must submit a collaboration log.

%% ============================================================
\section*{Part A: Event Study Basics (25 points)}
%% ============================================================

We consider an event study approach using data from Sun and Abraham (2021), which replicates results from Dobkin et al.\ (2020) using the HRS data. Variables: \texttt{hhidpn} (household identifier---the individual panel ID), \texttt{wave} (time identifier), \texttt{wave\_hosp} (time of event---when the individual is hospitalized), and \texttt{oop\_spend} (out-of-pocket spending).

We follow Sun and Abraham's notation throughout. Denote the initial time period of treatment for unit $i$ as $E_i$.

\begin{enumerate}

    \item \textbf{Treatment indicator.} What variable corresponds to $E_i$ in the dataset? Construct $D_{it} = \mathbf{1}(E_i \leq t)$, equal to one when an individual is treated. Report the share of individuals treated in periods 7, 8, 9, and 10.


    \item \textbf{Static TWFE.} Estimate the traditional static two-way fixed effects specification:
    \begin{equation}\label{eq:twfe}
        Y_{it} = \alpha_i + \lambda_t + D_{it}\beta + \epsilon_{it}
    \end{equation}
    where $Y_{it}$ is \texttt{oop\_spend}, $\alpha_i$ is a unit fixed effect, and $\lambda_t$ is a time fixed effect. Report the estimate for $\beta$ and its standard error (adjust for appropriate inference in the panel setting).


    \item \textbf{Group-by-group estimation.} Denote the control group as the last cohort ever treated. For each other treated wave, estimate the treatment effect relative to this group, excluding the last period of data. Report the coefficients and standard errors for each wave.

    How do these estimates compare to the static TWFE result? For the Wave 8 cohort, the diff across units is Cohort Wave 8 vs.\ Cohort Wave 11. What is the diff across time comparing?


    \item \textbf{Dynamic TWFE.} Define $D_{it}^l = \mathbf{1}(t - E_i = l)$ and estimate:
    \begin{equation}\label{eq:twfe_dyn}
        Y_{it} = \alpha_i + \lambda_t + \sum_{l \in \{-3,-2\}} D_{it}^l \beta_l + \sum_{l \in \{0,1,2,3\}} D_{it}^l \beta_l + \epsilon_{it}
    \end{equation}
    Report the $\beta$ coefficients and their standard errors.


\end{enumerate}

%% ============================================================
\section*{Part B: Event Study Extensions (25 points)}
%% ============================================================

\begin{enumerate}[resume]

    \item \textbf{Cohort-specific dynamic effects.} Repeat the dynamic specification from Equation~\ref{eq:twfe_dyn}, but estimate group-by-group. Focusing on Cohort Wave 8 vs.\ Cohort Wave 11: How would you run this specification? What coefficients can you estimate? Report these estimates.

    Now estimate $\beta_0$ for each cohort separately. How do these estimates compare to Equation~\ref{eq:twfe_dyn}?


    \item \textbf{Decomposing the pre-trend test.} Focus on $\beta_{-2}$ from Equation~\ref{eq:twfe_dyn}---the traditional pre-trend test. Sun and Abraham show that under standard diff-in-diff assumptions:
    \begin{equation}
        \beta_{-2} = \sum_{e=8}^{11} \omega_{e,-2}^{-2} CATT_{e,-2} + \sum_{l \in \{-3,0,1,2,3\}} \sum_{e=8}^{11} \omega_{e,l}^{-2} CATT_{e,l} + \sum_{l' \in \{-4,-1\}} \sum_{e=8}^{11} \omega_{e,l'}^{-2} CATT_{e,l'}
    \end{equation}
    where $CATT_{e,l}$ is the average treatment effect $l$ periods from initial treatment for cohort $e$. Estimate the weights $\omega$ by replacing $Y_{it}$ in Equation~\ref{eq:twfe_dyn} with $D_{it}^l \mathbf{1}(E_i = e)$ as the outcome and reporting the coefficient on $D_{it}^{-2}$. Do so for each $l$ and $e$.

    Your results should match Figure 2 in Sun and Abraham. How does this affect your interpretation of the pre-trend test?

    \item \textbf{Sun and Abraham estimator.} Estimate the interaction-weighted estimator that avoids contamination bias:
    \begin{equation}
        Y_{it} = \alpha_i + \lambda_t + \sum_{e \in \{8,9,10\}} \sum_{\substack{l = -3 \\ l \neq -1}}^{3} \mathbf{1}(E_i = e) \times D_{it}^l \, \delta_{e,l} + \epsilon_{it}
    \end{equation}
    excluding the last time period and treating Cohort Wave 11 as the control group. Report $\delta_{e,0}$ for all three treated cohorts.

    The final estimate $\mu_0$ weights each $\delta_{e,0}$ by the cohort sample weight $\pi_e = \Pr(E_i = e \mid l = 0)$. Report this estimate.


\end{enumerate}

%% ============================================================
\section*{Part C: Synthetic Methods I (25 points)}
%% ============================================================

We now implement synthetic control approaches. You should use the \texttt{synthdid} package (\url{https://synth-inference.github.io/synthdid/}) and the \texttt{scpi} package (\url{https://nppackages.github.io/scpi}).

\begin{enumerate}[resume]

    \item \textbf{Synthetic control.} Replicate the main synthetic control effect of California's Proposition 99 on cigarette consumption. Set up the data following the \texttt{synthdid} package vignette. Report the treatment effect for the synthetic control estimator (note: this is \emph{not} the main \texttt{synthdid} estimator).


    \item \textbf{Synthetic DID.} Estimate the effect using the \texttt{synthdid} estimator. Report the treatment effect and compare to the synthetic control effect.


    \item \textbf{Comparing weights.} Examine the weights for the synthetic control. What units are most important? How do these weights compare to the weights in the synthetic DID estimator? (In R, use the \texttt{summary} command on the estimate object.)

    \item \textbf{Alternative SC implementation.} Use the \texttt{scpi} package from Cattaneo et al.\ to estimate the effect of California's Prop 99. See the illustration at \url{https://github.com/nppackages/scpi/blob/main/R/scpi_illustration.R}. Implement the simplex version of SC.

    Hint: the estimate object includes \texttt{est.results\$Y.post.fit} with year-by-year synthetic values. Contrast these with actual California values in the post period and take the overall mean difference.

    Compare this with the \texttt{synthdid} synthetic control estimate (they should be close but likely not identical).


    \item \textbf{Prediction intervals.} Estimate prediction intervals from Cattaneo et al.\ using default settings. Report the synthetic California estimate and the left and right bounds for the year 2000. Is the true value inside that interval?


    \item \textbf{Placebo test.} Pretend the treatment was in 1985 and rerun the estimation procedure. Prior to 1989, is synthetic California significantly different from actual California? (You can plot this using \texttt{sc\_plot}.)

    Report whether actual California is inside the confidence intervals for synthetic California as of 1988.


    \item \textbf{Placebo comparison.} Report the synthetic California estimate and bounds in the year 2000 under the placebo treatment timing.


\end{enumerate}

%% ============================================================
\section*{Part D: Synthetic Methods II --- Staggered Design (25 points)}
%% ============================================================

We now replicate results from Figure 3 and Table 2 of Wolfers (2006) using \texttt{wolfers\_aer\_clean.csv}. This dataset has been cleaned for panel balance (48 states, 1959--1990). We focus on estimating post-treatment effects only. Use the \texttt{synthdid} package.

Variables: \texttt{st} (state), \texttt{year}, \texttt{div\_rate} (divorce rate, outcome), \texttt{unilateral} (treatment indicator), \texttt{lfdivlaw} (year of divorce law; 2000 if never treated).

\begin{enumerate}[resume]

    \item \textbf{Treatment timing.} Load the dataset and create a balanced panel. Report the distribution of treatment timings. What two years have the largest share of divorce laws passed?


    \item \textbf{Single-cohort synthdid.} The \texttt{synthdid} package is not designed for staggered settings. Using only states treated in 1971 or never treated, estimate the overall synthdid effect with the 1971 cohort as the treated group.


    \item \textbf{Dynamic effects for 1971 cohort.} Using the same sample (1971-treated and never-treated states), estimate dynamic effects following Wolfers (2006). Estimate the effect in years 1--2 (1971--1972), years 3--4 (1973--1974), and so on up to 15+ years after, pooling all remaining years.

    Report the estimated effects and compare to Table 2 of Wolfers. How do the point estimates compare? If you implement standard errors, do Wolfers' estimates lie within the 95\% confidence interval?


    \item \textbf{Aggregating across cohorts.} Implement the cohort-by-cohort approach for each treatment year (always using untreated states as controls). Combine estimates into a single estimate using the relative share of each treatment timing as weights. Ignore the always-treated state.

    How does this compare to Table 2 of Wolfers?


\end{enumerate}

\begin{boxK}
\section*{Conceptual Questions} \textbf{AI Free Zone}

\begin{enumerate}[resume]

    \item \textbf{TWFE comparisons.} Think back to the traditional static TWFE equation (Equation~\ref{eq:twfe}). What is the relative comparison period for this diff-in-diff? How does this differ from the group-by-group comparisons in Part A, and why does this matter for interpretation?

    \item \textbf{Interpreting Wolfers.} Wolfers (2006) concludes: ``I find that the divorce rate rose sharply following the adoption of unilateral divorce laws, but that this rise was reversed within about a decade. There is no evidence that this rise in divorce is persistent.'' Are your synthetic control results consistent with this conclusion? In what ways do they agree or disagree, and what might explain any differences?

\end{enumerate}
\end{boxK}


\section*{Submission Requirements}

Submit to GitHub Classroom:
\begin{itemize}
    \item \texttt{homework6-code.R} --- your code
    \item \texttt{homework6-writeup.pdf} --- answers to written questions
    \item \texttt{homework6-collab.txt} --- collaboration log
\end{itemize}

\paragraph{Collaboration log:} List any AI tools used. For each, briefly note what you used it for and whether you had to correct its output. If you didn't use AI tools, write ``No AI tools used.''

\end{document}
